前面的求解器保持网络不变，主要减少积分所需的调用次数。
蒸馏则重新训练学生，使一次学生调用完成教师多次调用完成的事情。
例如教师用两个小步从 $x_t$ 到达 $x_s$，记结果为
$\hat x_s=\Phi_{r\to s}(\Phi_{t\to r}(x_t))$。
学生直接预测这两个时刻间的映射 $G_\theta(x_t,t,s)$，可最小化
\[
\mathbb E\|G_\theta(x_t,t,s)-\operatorname{sg}(\hat x_s)\|^2.
\]
$\operatorname{sg}$ 表示目标不参与反向传播。
逐次将两个步骤合为一个，就是渐进式蒸馏的基本思路。
它学习的是教师给出的对应关系，教师本身的偏差和离线目标构造的成本也会随之带入。

另一条路线不要求学生逐样本模仿教师轨迹，而要求学生输出的整体分布接近教师分布。
设 $x_0=G_\theta(z)$，在噪声级 $t$ 上加噪，学生诱导的边缘分布记为 $q_t^\theta$。
分布匹配的一类理想目标是
$\mathbb E_t[w(t)D_{\rm KL}(q_t^\theta\|p_t^{\rm teacher})]$。
在可交换求导等条件下，其梯度涉及
\[
\mathbb E\left[w(t)a_t
 \bigl(s_{q_t^\theta}(x_t)-s_{p_t^{\rm teacher}}(x_t)\bigr)^\top
 \partial_\theta G_\theta(z)\right].
\]
教师分数可以由预训练扩散模型提供，但学生分布的分数仍要估计。
学生不断变化，后一种估计也要跟着更新。
DMD2 等方法继续处理这种训练稳定性，并结合额外分布约束；
上式用于说明两种分布的分数为什么出现，不是完整复现它的全部训练目标。
与轨迹回归相比，这条路线允许学生改变噪声与输出的对应关系。
下面的一致性模型又从“同一条轨迹的不同点应有同一个终点”出发建立训练约束。
\cite{theorybook,dmd2}
